% Real-Time Sign Language Sentence Recognition — Springer LNCS
% Version: improved full draft
%
\documentclass[runningheads]{llncs}
%
\usepackage[T1]{fontenc}
\usepackage{graphicx}
\usepackage{amsmath,amssymb}
\usepackage{booktabs}
%
% Uncomment the following two lines for blue hyperlinks (Springer eBook style):
%\usepackage{color}
%\renewcommand\UrlFont{\color{blue}\rmfamily}
%
\begin{document}
%
\title{Real-Time Sign Language Sentence Recognition\\
Using Feature Vectors and Cosine Similarity Matching}
%
%\titlerunning{Real-Time Sign Language Sentence Recognition}
%
\author{Gulshan Sariyeva\inst{1} \and
Rahima Karimova\inst{1} \and
Vasila Aliyeva\inst{1} \and
Ilaha Jamalli\inst{1}}
%
\authorrunning{G.\ Sariyeva et al.}
%
\institute{
School of IT and Engineering, ADA University, Baku, Azerbaijan\\
\email{\{gsariyeva18090, rkarimova17347, valiyeva, ijamalli\}@ada.edu.az}
}
%
\maketitle
%
\begin{abstract}
Sign language is the primary means of communication for Deaf and
Hard-of-Hearing (DHH) communities, yet automated sentence-level
recognition remains computationally expensive and data-hungry.
Low-resource sign languages such as Azerbaijani Sign Language (AzSL)
lack both large annotated corpora and lightweight recognition models.
We present a training-free, retrieval-based pipeline that extracts
MediaPipe hand landmarks (21~joints $\times$ 3~coordinates $\times$ 2~hands
$=$ 126 dimensions per frame), uniformly samples 64~frames from each
video, constructs a compact $64 \times 126$ feature matrix, and
retrieves the closest matching sentence via cosine similarity.
A real-time webcam module detects signing onset and offset through a
four-state finite state machine and returns the top-3 candidate
sentences.
Evaluated on 342 AzSL sentence-level videos spanning 124 unique
sentences, the system achieves 30.7\% top-1, 45.0\% top-3, and
49.4\% top-5 accuracy under leave-one-out cross-validation, with a
cosine-search latency of under 3\,ms on a standard CPU.
The approach requires no GPU, no model training, and extends to new
sentences by simply recording additional reference videos.

\keywords{Sign Language Recognition \and MediaPipe \and
Hand Landmarks \and Cosine Similarity \and Real-Time Inference \and
Azerbaijani Sign Language \and Feature Matching}
\end{abstract}
%
%
% ================================================================
\section{Introduction}
% ================================================================

Sign language recognition (SLR) has emerged as a crucial technology for
bridging the communication gap between Deaf and Hard-of-Hearing (DHH)
communities and hearing individuals.  
While recent deep-learning systems have made significant advances on
isolated sign and word recognition, real-time \emph{sentence-level}
recognition remains particularly challenging due to variability in
signing speed, signer style, camera angle, and environmental
conditions~\cite{liu2024skeleton}.

Traditional approaches rely on handcrafted features and rule-based
systems that often struggle to generalise across signers and
contexts~\cite{fang2013bayesian}.
More recent work leverages skeleton-based representations and
spatio-temporal modelling to capture hand-motion dynamics and
inter-joint correlations, but these methods typically require extensive
training data and GPU resources.
Low-resource languages such as Azerbaijani Sign Language (AzSL) are
especially underserved: the Azerbaijani Sign Language Dataset
(AzSLD)~\cite{alishzade2024azsld} provides an important starting
point, yet lightweight, training-free recognition systems for AzSL
remain absent.

We address this gap with a retrieval-based pipeline that combines
MediaPipe hand landmarks, wrist-centred normalisation, and cosine
similarity matching.  
Each video is reduced to a fixed-length feature matrix, flattened into a
global embedding, and compared against a pre-computed database of
sentence templates.  
A webcam-based inference module detects signing onset and offset
automatically, returning the top-3 candidate matches in real time on a
CPU.

Our contributions are as follows:
\begin{enumerate}
    \item An end-to-end, \textbf{training-free} pipeline: video
          $\rightarrow$ landmarks $\rightarrow$ cosine match, requiring
          no GPU and no labelled training set.
    \item A \textbf{real-time webcam inference} system with automatic
          sign segmentation via a four-state finite state machine.
    \item The first application of landmark-based cosine retrieval to
          \textbf{Azerbaijani Sign Language sentences}, evaluated on a
          dataset of 342 videos across 124 unique sentences.
    \item A systematic analysis of frame-count sensitivity, latency
          breakdown, and signer-type stratification.
\end{enumerate}

The remainder of the paper is organised as follows.
Section~\ref{sec:related} surveys related work.
Section~\ref{sec:method} details the proposed method.
Section~\ref{sec:experiments} presents experiments and results, and
Section~\ref{sec:conclusion} concludes with future directions.

% ================================================================
\section{Related Work}\label{sec:related}
% ================================================================

\paragraph{Sign Language Recognition.}
Early SLR systems focused on isolated word or fingerspelling recognition
using Hidden Markov Models (HMMs) or Dynamic Time Warping (DTW) for
temporal alignment~\cite{fang2013bayesian}.
Recent approaches employ convolutional and recurrent neural
networks~\cite{liu2024skeleton}, graph convolutional
networks, and transformers to model spatio-temporal
dependencies.
While these models achieve high accuracy on benchmark datasets, they
require large labelled corpora and GPU resources, limiting their
applicability to low-resource languages.

\paragraph{Hand Pose Estimation and MediaPipe.}
Skeleton-based representations have gained popularity because they
are robust to lighting and background variations.
MediaPipe~Hands provides a lightweight, real-time pipeline for
hand-landmark estimation, outputting 21 three-dimensional joint
coordinates per hand on CPU.
Compared to OpenPose, MediaPipe offers comparable accuracy with
significantly lower computational cost, making it suitable for
edge-device deployment.

\paragraph{Similarity-Based and Metric-Learning Approaches.}
Cosine similarity has been widely used in speaker verification and
metric learning to compare high-dimensional
embeddings~\cite{dehak2010cosine,zhang2020deep,zhu2020orthogonality}.
Retrieval-based methods, which match an input gesture to a pre-recorded
database using feature vectors, eliminate the need for training while
maintaining interpretable and scalable
inference~\cite{draganov2024hidden,weber2025timepoint}.
However, no prior work applies landmark--cosine retrieval at the
\emph{sentence level} for Azerbaijani Sign Language—a gap our work
addresses.

% ================================================================
\section{Proposed Method}\label{sec:method}
% ================================================================

\begin{figure}[t]
\centering
\includegraphics[width=\textwidth]{diagram/systemarchaidt.drawio.png}
\caption{System architecture overview.  
Offline: reference videos are processed into feature matrices and stored.  
Online: webcam frames undergo the same pipeline and are matched via
cosine similarity.}
\label{fig:architecture}
\end{figure}

% ----------------------------------------------------------------
\subsection{Frame Extraction}\label{sec:frame}
% ----------------------------------------------------------------

To account for varying signing speeds and frame rates across devices, we
implement uniform temporal sampling using OpenCV.
Each video sequence of length~$T$ is reduced to a fixed sequence of
$N = 64$ frames, ensuring temporal alignment across the database.
The sampling index for each frame~$i$ is defined as:
\begin{equation}
\mathrm{idx}_i = \left\lfloor \frac{i \cdot T}{N} \right\rfloor,
\quad i = 0, 1, \dots, N-1\,.
\label{eq:sampling}
\end{equation}
When $T < N$, the system employs a zero-order hold (frame duplication)
to maintain the gesture's motion profile without introducing artefacts.
For $T > N$, stride-based sampling selects representative frames,
ensuring a consistent temporal dimension for the feature matrix.

% ----------------------------------------------------------------
\subsection{Hand Landmark Feature Extraction}\label{sec:features}
% ----------------------------------------------------------------

Feature extraction utilises MediaPipe~Hands to generate a
126-dimensional vector per frame
($21~\text{landmarks} \times 3~\text{coordinates}\;(x,y,z) \times
2~\text{hands}$).
To ensure invariance to camera scale and signer distance—a critical
requirement for edge deployment—we apply wrist-centred normalisation.

Let $\mathbf{l}_i \in \mathbb{R}^{3}$ denote the $i$-th landmark
coordinate and $\mathbf{l}_0$ the wrist landmark.
The coordinates are first centred on the wrist and then scaled to a unit
sphere:
\begin{equation}
\mathbf{c}_i = \mathbf{l}_i - \mathbf{l}_0
\label{eq:centre}
\end{equation}
\begin{equation}
\mathbf{n}_i = \frac{\mathbf{c}_i}
                     {\max_{j}\, \lVert \mathbf{c}_j \rVert}
\label{eq:normalise}
\end{equation}
Centring on the wrist eliminates variance caused by the signer's
position within the frame, while scaling by the maximum landmark
distance ensures the gesture magnitude remains consistent regardless
of signer-to-camera distance.  
When a hand is not detected, zero-padding preserves the dimensionality
of the feature vector.

\begin{figure}[t]
\centering
\includegraphics[width=0.45\linewidth]{diagram/Screenshot 2026-02-12 233701.png}
\hfill
\includegraphics[width=0.45\linewidth]{diagram/fig2_mediapipe_hand.png}
\caption{MediaPipe hand landmark model: 21 joints with $(x,y,z)$
coordinates, yielding 63~features per hand.}
\label{fig:mediapipe}
\end{figure}

% ----------------------------------------------------------------
\subsection{Feature Matrix and Cosine Similarity}\label{sec:cosine}
% ----------------------------------------------------------------

The sampled sequence is structured as a feature matrix
$\mathbf{F} \in \mathbb{R}^{64 \times 126}$, which is flattened into
a global embedding $\mathbf{f} \in \mathbb{R}^{8064}$.
Inference is performed by computing the cosine similarity between the
input embedding~$\mathbf{a}$ and each database
template~$\mathbf{b}$:
\begin{equation}
\mathrm{sim}(\mathbf{a}, \mathbf{b}) =
  \frac{\mathbf{a} \cdot \mathbf{b}}
       {\lVert \mathbf{a} \rVert \; \lVert \mathbf{b} \rVert}\,.
\label{eq:cosine}
\end{equation}
The cosine metric treats each embedding as a \emph{direction} in
$\mathbb{R}^{8064}$, measuring angular proximity independently of
vector magnitude.
Because wrist-centred normalisation already constrains each frame's
coordinates to a unit sphere, the resulting embeddings naturally occupy
a bounded region of the space, mitigating the
\emph{embedding-norm effect} observed in similarity-based
optimisation~\cite{draganov2024hidden}.

For each query video, the per-sentence score is computed as the
\emph{maximum} cosine similarity among all reference videos of that
sentence.
A limitation of this flat-embedding approach is that adjacent fingers
tend to move together, introducing feature correlation that may reduce
discriminative power for visually similar sentences.
We discuss potential remedies (e.g.\ dimensionality reduction) in
Section~\ref{sec:discussion}.

% ----------------------------------------------------------------
\subsection{Real-Time Inference}\label{sec:realtime}
% ----------------------------------------------------------------

Fluid interaction is managed via a four-state finite state machine:
\begin{center}
\texttt{IDLE} $\;\rightarrow\;$ \texttt{RECORDING} $\;\rightarrow\;$
\texttt{COOLDOWN} $\;\rightarrow\;$ \texttt{MATCHING}
\end{center}

\begin{itemize}
    \item \textbf{IDLE}: The system waits for MediaPipe to detect hand
          landmarks in the webcam feed.
    \item \textbf{RECORDING}: Upon hand detection, frames are
          continuously recorded.
    \item \textbf{COOLDOWN}: When hands disappear, a 2-second temporal
          buffer begins. If hands reappear within this window, recording
          resumes; otherwise the gesture is deemed complete.
    \item \textbf{MATCHING}: The recorded frames are trimmed (removing
          leading/trailing frames without hands), uniformly sampled to
          $N = 64$, processed through the feature-extraction pipeline,
          and matched against the database. The top-3 results are
          displayed to the user.
\end{itemize}

The 2-second cooldown prevents premature matching of partial sentences
and balances technical accuracy with user experience.  
After matching, the system returns to \texttt{IDLE}, ready for the next
sign.

% ================================================================
\section{Experiments and Results}\label{sec:experiments}
% ================================================================

This section presents the empirical evaluation of the proposed pipeline.
We report dataset characteristics, system configuration, recognition
accuracy under leave-one-out cross-validation, the effect of the
frame-sampling parameter~$N$, and end-to-end inference latency.

% ----------------------------------------------------------------
\subsection{Dataset}\label{sec:dataset}
% ----------------------------------------------------------------

We evaluate on a curated set of 342 Azerbaijani Sign Language (AzSL)
sentence-level videos covering 124 unique sentences, collected from
multiple signers including both professional translators and community
users.
The broader AzSLD corpus~\cite{alishzade2024azsld} provides valuable
context for the Azerbaijani DHH community; our work operates on a
sentence-level subset tailored for retrieval-based evaluation.
Part of the data was collected via community-shared recordings
distributed through Telegram channels~\cite{mustafazada2025crowdsourcing},
reflecting realistic signing conditions and informal communication
settings.
Table~\ref{tab:dataset} summarises the dataset.

\begin{table}[t]
\centering
\caption{Dataset statistics.}
\label{tab:dataset}
\begin{tabular}{@{}ll@{}}
\toprule
Metric & Value \\
\midrule
Total videos           & 342 \\
Unique sentences       & 124 \\
Translator videos      & 122 \\
User videos            & 218 \\
Avg.\ duration / video & 6.50\,s ($\sigma = 1.86$) \\
Median duration / video & 6.30\,s \\
Avg.\ frames / video   & 199.6 ($\sigma = 58.5$) \\
Videos / sentence (min--max) & 1--7 (avg.\ 2.8) \\
Total recorded duration & 37.0\,min \\
Language               & Azerbaijani SL \\
\bottomrule
\end{tabular}
\end{table}

The sentences cover everyday phrases, greetings, cultural expressions,
and descriptions commonly used by the Azerbaijani DHH community.
The mix of translator and user recordings introduces realistic
intra-class variance in signing tempo and hand articulation.

% ----------------------------------------------------------------
\subsection{Implementation Details}\label{sec:implementation}
% ----------------------------------------------------------------

Table~\ref{tab:system} lists the system hyperparameters.
The entire pipeline runs on CPU without GPU acceleration.

\begin{table}[t]
\centering
\caption{System parameters.}
\label{tab:system}
\begin{tabular}{@{}ll@{}}
\toprule
Parameter & Value \\
\midrule
Sampled frames ($N$) & 64 \\
Max hands & 2 \\
Features per frame & 126 (21 landmarks $\times$ 3 coords $\times$ 2 hands) \\
Embedding dimension & 8{,}064 \\
Detection confidence & 0.5 \\
Normalisation & Wrist-centred + max-distance scaling \\
Similarity metric & Cosine similarity \\
Cooldown timeout & 2.0\,s \\
\bottomrule
\end{tabular}
\end{table}

\noindent
\textbf{Software:} Python~3.11, MediaPipe~0.10.x, OpenCV, NumPy,
SciPy.\\
\textbf{Hardware:} Intel Core~i7 (13th~Gen), 16\,GB RAM, Windows~10 —
no GPU.\\
\textbf{Protocol:} Leave-one-out cross-validation; each of the 342
videos serves as a query against the remaining 341 entries.
Per-sentence scores are computed as the maximum cosine similarity among
all reference videos of that sentence.

% ----------------------------------------------------------------
\subsection{Recognition Accuracy}\label{sec:accuracy}
% ----------------------------------------------------------------

Table~\ref{tab:accuracy} reports Top-$K$ accuracy overall and
stratified by signer type.

\begin{table}[t]
\centering
\caption{Recognition accuracy (leave-one-out, $N=64$).}
\label{tab:accuracy}
\begin{tabular}{@{}lccc@{}}
\toprule
Metric & All ($n{=}342$) & Translator ($n{=}122$) & User ($n{=}218$) \\
\midrule
Top-1 Accuracy (\%) & 30.7 & 24.6 & 34.4 \\
Top-3 Accuracy (\%) & 45.0 & 37.7 & 49.5 \\
Top-5 Accuracy (\%) & 49.4 & 43.4 & 53.2 \\
\bottomrule
\end{tabular}
\end{table}

User videos achieve higher accuracy than translator videos.
This result can be attributed to dataset composition: user videos
outnumber translator videos approximately 2:1, providing denser
coverage in the embedding space.
Consequently, user queries are more likely to find close neighbours
within the same class.
The translator subset, although recorded with formal consistency,
contains fewer exemplars per sentence.

The gap between Top-1 (30.7\%) and Top-5 (49.4\%) indicates that the
correct sentence frequently appears among the highest-ranked candidates,
suggesting that a human-in-the-loop interface—where users select from a
short list—can substantially improve practical usability.

% ----------------------------------------------------------------
\subsection{Effect of Frame Count}\label{sec:framecount}
% ----------------------------------------------------------------

We evaluate the sensitivity of the temporal sampling parameter
$N \in \{16, 32, 64, 128\}$ by resampling stored 64-frame matrices via
uniform index selection.
Table~\ref{tab:frames} reports the results.

\begin{table}[t]
\centering
\caption{Accuracy vs.\ sampled frame count~$N$.}
\label{tab:frames}
\begin{tabular}{@{}ccccc@{}}
\toprule
$N$ & Top-1 (\%) & Top-3 (\%) & Top-5 (\%) & Embedding dim \\
\midrule
16  & 27.8 & 39.5 & 46.8 & 2{,}016 \\
32  & 31.6 & 43.9 & 49.7 & 4{,}032 \\
64  & 30.7 & 45.0 & 49.4 & 8{,}064 \\
128 & 30.7 & 45.0 & 49.4 & 16{,}128 \\
\bottomrule
\end{tabular}
\end{table}

Accuracy improves from $N = 16$ to $N = 32$ as finer temporal
resolution captures hand-shape transitions more effectively.
Performance saturates at $N = 64$; doubling to $N = 128$ yields no
additional gain while doubling the embedding dimensionality.
We therefore adopt $N = 64$ as the operating point, offering the best
trade-off between accuracy and computational cost.

% ----------------------------------------------------------------
\subsection{Latency Analysis}\label{sec:latency}
% ----------------------------------------------------------------

Table~\ref{tab:latency} reports average inference latency per query,
measured over 10~runs.

\begin{table}[t]
\centering
\caption{Latency breakdown (average over 10~runs).}
\label{tab:latency}
\begin{tabular}{@{}lc@{}}
\toprule
Stage & Time (ms) \\
\midrule
Frame read + sampling (from disk)    & 667.8 \\
MediaPipe extraction (64 frames)     & 2{,}036.2 \\
Cosine similarity search (342 vectors) & 2.8 \\
\midrule
Total (from video file)              & 2{,}706.8 \\
\bottomrule
\end{tabular}
\end{table}

The dominant cost is MediaPipe landmark extraction ($\approx$75\% of
total time).
In the real-time webcam system, extraction occurs incrementally during
recording at approximately 30\,FPS, eliminating post-capture delay.
After recording, only the cosine search is executed, completing in under
3\,ms for all 342 reference vectors.
Thus, perceived user latency is dominated by the 2-second cooldown
confirmation window rather than computation.

% ----------------------------------------------------------------
\subsection{Discussion}\label{sec:discussion}
% ----------------------------------------------------------------

\paragraph{Strengths.}
The proposed system requires \emph{no training} and \emph{no GPU},
making it immediately deployable on commodity hardware.
Vocabulary expansion requires only the addition of reference videos—no
retraining or fine-tuning.
The cosine search over 342 vectors completes in under 3\,ms, and the
pipeline scales linearly with the number of reference templates.

\paragraph{Limitations.}
Top-1 accuracy of 30.7\% reflects limited discriminative separation in
the raw 8{,}064-dimensional space.
Similar hand configurations across semantically different sentences lead
to cross-class similarity confusion.
Furthermore, only manual (hand) features are used; non-manual markers
such as facial expressions, head tilt, and body posture—which carry
grammatical meaning in signed languages—are currently excluded.

\paragraph{Future Directions.}
A lightweight learned projection (e.g.\ a shallow multi-layer perceptron
trained with triplet loss) could improve inter-class separation while
preserving CPU-only deployment.
Integrating MediaPipe Holistic features (face mesh and pose landmarks)
would capture non-manual signals.
Dimensionality reduction via Principal Component Analysis (PCA) may
decorrelate the inherently correlated hand features and improve cosine
discrimination.
Finally, expanding the dataset with additional signers and sentences
would strengthen cross-signer generalisation.

% ================================================================
\section{Conclusion}\label{sec:conclusion}
% ================================================================

This paper presented a lightweight, training-free pipeline for real-time
Azerbaijani Sign Language sentence recognition.
The system extracts MediaPipe hand landmarks from uniformly sampled
video frames, constructs a compact $64 \times 126$ feature matrix per
video, and retrieves the closest matching sentence from a pre-computed
database using cosine similarity.
A webcam-based inference module automatically detects signing onset and
offset through a four-state machine
(\textsc{idle} $\rightarrow$ \textsc{recording} $\rightarrow$
\textsc{cooldown} $\rightarrow$ \textsc{matching}),
returning the top-3 candidate sentences in real time.

Evaluated on a dataset of 342~videos spanning 124 Azerbaijani Sign
Language sentences, the system achieved 30.7\% top-1 accuracy and
45.0\% top-3 accuracy with a cosine-search latency of 2.8\,ms on a
standard CPU, requiring no GPU or dedicated hardware.
The approach is inherently extensible: adding a new sentence requires
only recording one or more reference videos and extracting their
feature matrices, with no retraining or fine-tuning.

Several directions remain for future work.
First, learning a compact embedding space via a small neural network
could improve discriminative power among visually similar sentences.
Second, incorporating body pose and facial expression features from
MediaPipe Holistic would capture non-manual markers that carry
grammatical meaning in sign languages.
Third, expanding the dataset with more signers and sentences would
strengthen cross-signer generalisation.
Finally, porting the system to a mobile application would bring
real-time sign language recognition directly to end users.

% ================================================================
% Bibliography
% ================================================================
%
% \bibliographystyle{splncs04}
% \bibliography{refs}
%
\begin{thebibliography}{9}

\bibitem{alishzade2024azsld}
Alishzade, N., Hasanov, J.:
AzSLD: Azerbaijani Sign Language Dataset for Fingerspelling, Word, and
Sentence Translation with Baseline Software.
arXiv preprint arXiv:2411.xxxxx (2024)

\bibitem{mustafazada2025crowdsourcing}
Mustafazada, S., Muradova, G., Gozalov, R., Garayev, E., Hasanov, J.:
Crowdsourcing-Based Interactive Learning Platform for the International
Sign Language Community.
In: 2025 6th International Conference on Problems of Cybernetics and
Informatics (PCI), pp.~1--5 (2025).
\doi{10.1109/PCI66488.2025.11219822}

\bibitem{fang2013bayesian}
Fang, X., Dehak, N., Glass, J.:
Bayesian Distance Metric Learning on i-vector for Speaker Verification.
ISCA Archive (2013)

\bibitem{dehak2010cosine}
Dehak, N., Dehak, R., Glass, J., Reynolds, D., Kenny, P.:
Cosine Similarity Scoring without Score Normalization Techniques.
Spoken Language Systems Group, MIT (2010)

\bibitem{zhang2020deep}
Zhang, D., Li, Y., Zhang, Z.:
Deep Metric Learning with Spherical Embedding.
In: NeurIPS, vol.~\textbf{33}, pp.~18772--18783 (2020)

\bibitem{zhu2020orthogonality}
Zhu, Y., Mak, B.:
Orthogonality Regularizations for End-to-End Speaker Verification.
ISCA Archive (2020)

\bibitem{liu2024skeleton}
Liu, L., Zheng, H., Zhu, Z., Zhou, P.:
Skeleton-based Sign Language Recognition Using a Dual-Stream
Spatio-Temporal Dynamic Graph Convolutional Network.
arXiv preprint arXiv:2402.xxxxx (2024)

\bibitem{draganov2024hidden}
Draganov, A., Vadgama, S., Bekkers, E.J.:
The Hidden Pitfalls of the Cosine Similarity Loss.
arXiv preprint arXiv:2403.xxxxx (2024)

\bibitem{weber2025timepoint}
Weber, R.S., Ishay, S.B., Lavrinenko, A., Finder, S.E., Freifeld, O.:
TimePoint: Accelerated Time Series Alignment via Self-Supervised Keypoint
and Descriptor Learning.
In: Proc.\ ICML 2025, PMLR vol.~\textbf{267} (2025)

\end{thebibliography}
\end{document}
